\documentclass[a4paper]{exam}
\usepackage[svgnames]{xcolor}
\usepackage{amsmath,amssymb,amsthm,titlesec,minted}

\newcommand
\lralign[2]{%
  \noindent\makebox[\textwidth]{\(\displaystyle#1\)\hfill\ifx&#2&\else(#2)\fi}\vspace{2ex}
}

\renewcommand{\thesubsection}{\alph{subsection}.}
\renewcommand{\thesubsubsection}{\roman{subsubsection}.}

\definecolor{faintgreen}{rgb}{0.5,1,0.5}
\definecolor{faintred}{rgb}{1,0.5,0.5}
\definecolor{faintblue}{rgb}{0.5,0.5,1}
\definecolor{faintyellow}{rgb}{1,1,0.5}
\definecolor{codebg}{RGB}{40,42,54}

\titleformat{\section}{\normalfont\Large\bfseries}{}{0pt}{}
\titleformat{\subsection}{\normalfont\large\bfseries}{}{0pt}{}
\titleformat{\subsubsection}{\normalfont\normalsize\bfseries\itshape}{}{0pt}{}
\titlespacing*{\section}{1em}{1.6\baselineskip}{0.6\baselineskip}
\titlespacing*{\subsection}{2em}{1.2\baselineskip}{0.4\baselineskip}
\titlespacing*{\subsubsection}{4em}{1.0\baselineskip}{0.3\baselineskip}

\setminted{style=dracula,bgcolor=codebg,breaklines,autogobble,fontsize=\small,tabsize=4,baselinestretch=1}
\NewDocumentCommand{\codefromfile}{O{}mm}{\par\vspace{0.4em}\inputminted[#1]{#2}{#3}\par\vspace{0.4em}}
\NewDocumentEnvironment{Section}{m}
{%
  \section*{#1}%
  \begingroup
  \leftskip=1em
}
{%
  \par
  \endgroup
}

\NewDocumentEnvironment{Subsection}{m}
{%
  \subsection*{#1}%
  \begingroup
  \leftskip=2em
}
{%
  \par
  \endgroup
}

\setlength{\parindent}{0in}
\setlength{\oddsidemargin}{0in}
\setlength{\textwidth}{6.5in}
\setlength{\textheight}{8.8in}
\setlength{\topmargin}{0in}
\setlength{\headheight}{18pt}

\printanswers
\title{Week 1}
\author{Radhesh Goel}
\date{\today}

\begin{document}
\maketitle

% Question 1: Natural to Binary Numbers
\section
*{Question 1.}
\begin{solution}
  Divide repeatedly by \(2\), keeping track of the remainders and then reading
  remainders bottoms-up.
  \begin{Section}{1.}
    %-% 1. Compute the binary representation of the number 17.
    \[
      \begin{aligned}
        \frac{17}{2} & = 8 \text{ rem } 1 \\
        \frac{8}{2}  & = 4 \text{ rem } 0 \\
        \frac{4}{2}  & = 2 \text{ rem } 0 \\
        \frac{2}{2}  & = 1 \text{ rem } 0 \\
        \frac{1}{2}  & = 0 \text{ rem } 1
      \end{aligned}
    \]
    It gives us, \(\boxed{17_{10}=(10001)_2}\)
  \end{Section}

  \begin{Section}{2.}
    %-% 2. Compute the binary representation of the number 227.
    \[
      \begin{aligned}
        \frac{227}{2} & = 113 \text{ rem } 1 \\
        \frac{113}{2} & = 56 \text{ rem } 1 \\
        \frac{56}{2}  & = 28 \text{ rem } 0 \\
        \frac{28}{2}  & = 14 \text{ rem } 0 \\
        \frac{14}{2}  & = 7 \text{ rem } 0 \\
        \frac{7}{2}   & = 3 \text{ rem } 1 \\
        \frac{3}{2}   & = 1 \text{ rem } 1 \\
        \frac{1}{2}   & = 0 \text{ rem } 1
      \end{aligned}
    \]
    Therefore, \(\boxed{227_{10}=(11100011)_2}\)
  \end{Section}
  \begin{Section}{3.}
    %-% 3. Compute the binary representation of the number 55.4
    First, convert the whole number part:
    \[
      \begin{aligned}
        \frac{55}{2} & = 27 \text{ rem } 1 \\
        \frac{27}{2} & = 13 \text{ rem } 1 \\
        \frac{13}{2} & = 6 \text{ rem } 1 \\
        \frac{6}{2}  & = 3 \text{ rem } 0 \\
        \frac{3}{2}  & = 1 \text{ rem } 1 \\
        \frac{1}{2}  & = 0 \text{ rem } 1
      \end{aligned}
    \]
    Hence, \(\boxed{55_{10}=(110111)_2}\) For the fractional part:
    \[
      \begin{aligned}
        0.4 \times 2 & = 0.8 &  & \Rightarrow 0 \\
        0.8 \times 2 & = 1.6 &  & \Rightarrow 1 \\
        0.6 \times 2 & = 1.2 &  & \Rightarrow 1 \\
        0.2 \times 2 & = 0.4 &  & \Rightarrow 0
      \end{aligned}
    \]
    which gives, \(0.4_{10} = (0.\overline{0110})_2\) \newline\newline
    Therefore, \(\boxed{55.4_{10}=(110111.\overline{0110})_2}\)
  \end{Section}
\end{solution}

% Question 2: Binary to natural numbers
\section
*{Question 2.}
\begin{solution}
  Each digit is multiplied by its corresponding power of \(2\).
  \begin{Section}{1.}
    %-% 1. Determine the natural number value of (1010101)_2.
    \[
      \begin{aligned}
        (1010101)_2 & = 1(2^6)+0(2^5)+1(2^4)+0(2^3)+1(2^2)+0(2^1)+1(2^0) \\
                    & = 64 + 16 + 4 + 1 \\
                    & = 85
      \end{aligned}
    \]
    Hence, \(\boxed{(1010101)_2=85_{10}}\)
  \end{Section}
  \begin{Section}{2.}
    %-% 2.  Determine the natural number value of (11.1010...)_2.
    For whole number part,
    \[
      (11)_2=1(2^1)+1(2^0)=3
    \]
    For fractional part,
    \[
      (0.101010\ldots)_2 = \frac{1}{2} + \frac{1}{2^3} + \frac{1}{2^5} + \cdots
    \]
    The sum of this GP is \(\frac{2}{3}\) \newline\newline Therefore,
    \[
      \begin{aligned}
        (11.\overline{10})_2 & = 3 + \frac{2}{3} \\
                             & = \frac{11}{3} \\
                             & = (3.\overline{6})_{10}
      \end{aligned}
    \]
    Hence, \(\boxed{(11.\overline{10})_2 = 3.\overline{6}_{10}}\)
  \end{Section}
\end{solution}

% Question 3: General short response
\section
*{Question 3.}
\begin{solution}

  \begin{Section}{1.}
    %-% 1. In the IEEE754 single-precision (32-bit) floating point standard, how are the 32-bits allocated
    %-% between the sign, exponential and mantissa?
    The 32 bits are divided as follows:
    \[
      \boxed{1\text{ sign bit} \;+\; 8\text{ exponent bits} \;+\; 23\text{ mantissa bits}}
    \]
    The sign bit determines whether the number is positive or negative. The
    exponent controls the scale of the number, while the mantissa stores its
    significant digits. For normalised values, the leading \(1\) is implied,
    giving effectively \(24\) bits of precision.
  \end{Section}
  \begin{Section}{2.}
    %-% 2. A mechanical engineer is running an iterative stress-analysis calculation in single-precision float-
    %-% ing point. Machine epsilon (εm) for IEEE 754 single precision is approximately εm = 1.19×10−7.
    %-% What does this value represent, and why does it matter for the engineer’
    Machine epsilon,
    \[
      \varepsilon_m \approx 1.19 \times 10^{-7}
    \]
    is the smallest value such that
    \[
      1 + \varepsilon_m>1
    \]
    can be represented as a different single-precision floating-point number. It
    means single precision provides roughly \(7\) significant decimal digits.
    \newline\newline This matters because each calculation may introduce a small
    rounding error. During many stress-analysis iterations, these errors can
    accumulate or affect whether the solution appears to converge. The engineer
    should therefore avoid using a convergence tolerance smaller than the
    available floating-point precision.
  \end{Section}

  \begin{Section}{3.}
    The most likely explanation is that decimal values such as \(0.3\) cannot be
    represented exactly in binary floating-point. The two formulas may therefore
    produce slightly different stored values due to rounding, even though both
    are displayed as \(0.3000\). \newline\newline

    For example, \(a = 0.30000000000000004, \qquad b = 0.29999999999999999\)

    Since \texttt{a == b} compares the exact stored values, the condition
    evaluates as false. A tolerance should be used instead:
    \[
      |a - b| < \varepsilon
    \]
  \end{Section}
\end{solution}

\end{document}
